\chapter{Assessment and skills acquired}\label{ch:assessment}

This chapter steps back from the three projects to draw out the skills acquired and the common methodological pattern.

\section{A common pattern: diagnose, decouple, harden}

The three projects share the same outline, summarised in Table~\ref{tab:synthesis}. The first two are migrations: the true source of fragility was not the business logic but a shared infrastructure dependency, and the right answer was to \emph{decouple} before \emph{hardening}, without ever interrupting the service. The third, regulatory in nature, applies the same rigour to a brand-new chain.

\begin{table}[H]
\centering\small
\renewcommand{\arraystretch}{1.35}
\begin{tabularx}{\textwidth}{@{}p{2.6cm}XX@{}}
\toprule
\textbf{Axis} & \textbf{Notion project} & \textbf{Tableau project} \\
\midrule
Triggering constraint & \texttt{notion-client} 0.9$\rightarrow$3.0 in the base image: breaking API + silent truncation & Tableau client too old in the base image: JWT blocked \\
Legacy fragility & No truncation detection & 14 manually rotated PATs + token pool + XCom hack \\
% Decoupling & To \texttt{qonto-python-dlt} (dlt Snowflake destination) & To \texttt{qonto-python-data-engineering} (K8s pods) \\
Real difficulty & Byte-for-byte parity + zero-downtime cutover & Unblocking the JWT by moving the dependency out + avoiding pod starvation \\
Resilience gained & Native pagination, load auditing. & A single secret, no more rotation \\
\bottomrule
\end{tabularx}
\caption{Summary of the pattern common to the two migrations.}
\label{tab:synthesis}
\end{table}

\section{Technical skills}

This internship let me consolidate and acquire a set of production-grade data-engineering skills:

\begin{itemize}
  \item \textbf{Ingestion and warehousing}: command of the EL paradigm, of the dlt framework (resources, pipelines, column hints, write dispositions), and of Snowflake loading mechanisms (S3 staging, \texttt{COPY INTO}, atomic swap, grant handling).
  \item \textbf{Orchestration}: designing Airflow DAGs, decoupling patterns via Kubernetes pods, sensors and dependency management, concurrency pools.
  \item \textbf{Security and authentication}: implementing JWT authentication (Connected App), understanding the PAT, OAuth and mTLS/PKI models, and secrets management through an encrypted parameter store.
  \item \textbf{Software quality}: per-module unit tests, continuous integration, linting, configuration validation, and tooling to automate recurring tasks.
  \item \textbf{Systems integration}: transport protocols (FTPS/mTLS), XML formats and schema validation.
\end{itemize}

\section{Methodological skills}

Beyond the technical side, the internship reinforced a way of working: \emph{trace problems back to the root cause} before coding (both migrations were born from a refusal to treat a symptom); \emph{make trade-offs explicit} rather than suffer them (choice of dlt, the Tableau server limit); \emph{design for reversibility and zero interruption}; and \emph{measure} the effect of changes (tasks reduced, an incident class eliminated). The team's \emph{Value Analysis / Value Engineering} approach served as a guiding thread for each of these points.

\section{Transferable skills}

I gained autonomy and \emph{ownership}: taking a high-impact topic and carrying it into production entails real responsibility, especially for critical pipelines that page the on-call rotation. Collaboration was essential: demanding code reviews, exchanges with different teams by documenting decisions, sticking points and trade-offs on Notion, both to transfer ownership of the systems to the team.

\section{Value for the company}

The deliverables carry direct and lasting value: the removal of two incident classes (silent truncation; missed token rotation), a significant reduction of legacy code and of coupling to the shared image, and the establishment of a non-existent regulatory chain. Beyond that, the patterns established (dlt's first Snowflake destination, a zero-downtime migration template, reusable JWT authentication) are references for the team's future migrations.
